\section*{Hoofdstuk \ref{ch:b3:conservation-biology} --- Conservatiebiologie}

\begin{solution}{exo:b3:conservation-biology:1}
\omterm{def:b3:conservation-biology:small}{Demografische toevalligheid}: toeval in de geboorten, de sterfgevallen
en de geslachten van een handvol individuen --- de kakapo bij
eenenvijftig, met meer mannetjes dan vrouwtjes. \omterm{def:b3:conservation-biology:small}{Toevalligheid van de
omgeving}: een slecht jaar dat elk individu tegelijk treft --- een
droogte, een ziekte, een harde winter die een populatie van een paar
honderd treft. \omterm{def:b3:conservation-biology:small}{Allee-effect}: een groei per hoofd die bij lage dichtheid
daalt --- de trekduif, wier kolonies niet meer konden broeden zodra zij
waren uitgedund, en wier overgebleven roofdieren niet meer werden
verzadigd.
\end{solution}

\begin{solution}{exo:b3:conservation-biology:2}
De omvang van een ideale populatie --- constant, willekeurig parend,
met gelijke geslachten en Poisson-verdeelde gezinsgrootten --- die met
hetzelfde tempo heterozygotie zou verliezen als de echte. Ongelijke
aantallen broedende mannetjes en vrouwtjes; schommelingen in de omvang,
die de kleinste generaties zwaarder laten wegen; een variantie in de
gezinsgrootte boven Poisson, waarbij enkele individuen de meeste jongen
voortbrengen (en ook niet-willekeurige paring en overlappende
generaties).
\end{solution}

\begin{solution}{exo:b3:conservation-biology:3}
$S = cA^{z}$. Het behouden deel $= (A'/A)^{z}$: $0.5^{0.25} = 0.84$;
$0.1^{0.25} = 0.56$; $0.01^{0.25} = 0.32$ --- een honderdste van de
oppervlakte behoudt uiteindelijk nog altijd een derde van de soorten.
\end{solution}

\begin{solution}{exo:b3:conservation-biology:4}
$N_{e} \ge 50$ houdt de aanwas van de inteelt onder $1/100$ per
generatie, het tempo waarbij dierfokkers de depressie verdragen: dat
beschermt tegen verlies van fitness op korte termijn. $N_{e} \ge
500$ brengt het verlies aan variatie door drift in evenwicht met de
aanvoer door mutatie, zodat de populatie de kwantitatieve variatie
behoudt die zij nodig heeft om zich aan te passen: dat beschermt de
lange termijn. Echte populaties hebben vijf- tot tienmaal deze getallen
in de telling nodig.
\end{solution}

\begin{solution}{exo:b3:conservation-biology:5}
$N_{e} = 4\times 8\times 120/128 = 30$; de heterozygotie daalt per
generatie met $1/60 =
\qty{1.7}{\%}$; een kwart verloren wanneer $(1 - 1/60)^{t} =
0.75$, dus $t = \ln 0.75/\ln(59/60) = 17$ generaties.
\end{solution}

\begin{solution}{exo:b3:conservation-biology:6}
$1/N_{e} = \tfrac{1}{5}(4/1000 + 1/20) = 0.0108$, $N_{e} = 93$: één
generatie bij twintig heeft vijf generaties zich als drieënnegentig
laten gedragen. Behouden heterozygotie $(1 - 1/186)^{5} = 0.973$, tegen
$(1 -
1/2000)^{5} = 0.9975$ bij een constante duizend.
\end{solution}

\begin{solution}{exo:b3:conservation-biology:7}
$F = 1 - (1 - 1/20)^{5} = 0.23$; fitness omlaag met $5\times 2.3 = \qty{11}{\%}$.
Elke nakomeling van een Texaans vrouwtje en een Floridaans mannetje
heeft $F = 0$; als de acht nieuwkomers twee vijfde van de broedende
vrouwtjes zijn, is twee vijfde van de volgende generatie uitgekruist en
daalt de gemiddelde $F$ in één generatie tot ongeveer
$0.6\times 0.23 = 0.14$ --- wat het herstel liet zien.
\end{solution}

\begin{solution}{exo:b3:conservation-biology:8}
$r = \ln(100/31)/8 = \qty{0.15}{yr^{-1}}$. Doorgezet tot 2020:
$100\,\mathrm{e}^{0.15\times 17} \approx 1200$. Het park herbergde er
ongeveer honderd: de territoria raakten vol, de wapiti's liepen terug,
roedels doodden elkaars leden, en schurft en hondenziekte verspreidden
zich --- de dichtheidsafhankelijkheid zette in zodra het lege gebied
bezet was.
\end{solution}

\begin{solution}{exo:b3:conservation-biology:9}
(a) Een achteruitgang van \qty{60}{\%} in tien jaar overtreft
\qty{50}{\%}: Bedreigd. (b) $180 < 250$ volwassen exemplaren: Ernstig
Bedreigd. (c) Verspreidingsgebied \qty{3000}{km^2}, onder \qty{5000}{}:
Bedreigd. (d) \num{30000} exemplaren en een achteruitgang van
\qty{10}{\%}: geen van beide drempels gehaald --- niet bedreigd
(hooguit Gevoelig, in de gaten te houden).
\end{solution}

\begin{solution}{exo:b3:conservation-biology:10}
Alle $N$ paren falen met kans $(1/2)^{N}$: $0.25$, $0.031$, $0.001$,
$10^{-6}$. Het demografische risico daalt exponentieel met $N$ en is
voorbij een paar dozijn individuen verwaarloosbaar; daarboven is het
risico dat overblijft dat van de omgeving, dat niet zo snel daalt,
omdat het alle individuen tegelijk treft.
\end{solution}

\begin{solution}{exo:b3:conservation-biology:11}
Eén snipper van een tiende: $0.1^{0.25} = 0.56$ van de oorspronkelijke
soorten. Tien snippers herbergen tussen $0.56$ (als zij alle dezelfde
soorten herbergen) en, in beginsel, meer dan het geheel (als elk andere
soorten herbergde --- onmogelijk boven de regionale voorraad). Waar zij
uitkomen hangt af van het verloop tussen de snippers: verschillende
leefgebieden in verschillende snippers duwen het totaal omhoog; soorten
die grote oppervlakten of binnenkant nodig hebben, of die door dezelfde
\omterm{prop:b3:conservation-biology:area}{randeffecten} uit elke snipper verdwijnen, duwen het omlaag; en de
uitstervingsschuld van elke snipper wordt afzonderlijk betaald, zodat
het totaal op lange termijn gewoonlijk onder dat van het geheel ligt.
\end{solution}

\begin{solution}{exo:b3:conservation-biology:12}
Drift neemt heterozygotie weg met tempo $1/2N_{e}$ per generatie;
immigranten vervangen een deel $m$ van de genenpoel door allelen die
elders zijn getrokken. Bij $mN_{e} \approx 1$ is $m \approx 1/N_{e}$,
tweemaal het tempo van het verlies, zodat de variatie sneller wordt
aangevuld dan zij vervalt en de populatie in allelfrequentie dicht bij
de bron blijft ($F_{ST} = 1/(1
+ 4N_{e}m) \approx 0.2$). De kosten: plaatselijk begunstigde allelen
worden per generatie met $m$ verdund, zodat de plaatselijke aanpassing
alleen overleeft voor kenmerken wier selectiecoëfficiënt ongeveer
$1/N_{e}$ overtreft; zwak geselecteerde plaatselijke verschillen worden
overspoeld.
\end{solution}

\begin{solution}{pb:b3:conservation-biology:1}
\textbf{1.} $N_{e} = 4\times 20\times 40/60 = 53$, tegen een telling van 60.
\textbf{2.} $1/N_{e} = \tfrac{1}{5}(1/500 + 1/100 + 1/60 + 1/51 + 1/60) =
0.0130$, $N_{e} = 77$; de flessenhalsgeneraties van 51 en 60
overheersen, de 500 telt nauwelijks mee.
\textbf{3.} $F = 1 - (1 - 1/154)^{5} = 0.032$; lineair $5/154 = 0.032$.
\textbf{4.} $6\times 0.32 = \qty{1.9}{\%}$.
\textbf{5.} $(1 - 1/154)^{5} = 0.968$: \qty{96.8}{\%} behouden.
\textbf{6.} Tien generaties bij $N_{e} = 53$: aanwas $1 - (1 -
1/107)^{10} = 0.090$, totaal $F = 1 - 0.968\times 0.910 = 0.12$;
verlies aan fitness $6\times 1.2 = \qty{7}{\%}$; een eeuw.
\textbf{7.} Paren van twee inheemse vogels vormen $(60/70)^{2} = 0.73$
van de paringen en hun nakomelingen houden $F \approx 0.03$ (plus een
kleine aanwas); alle andere hebben $F = 0$: gemiddelde
$F \approx 0.73\times 0.035 = 0.026$, een daling met een kwart in één
generatie.
\textbf{8.} $F$ is identiteit door afstamming binnen een individu, en
elke uitgekruiste nakomeling zet haar op nul terug; de heterozygotie is
een eigenschap van allelfrequenties, en de allelen van de migranten
verspreiden zich alleen naarmate zij worden doorgegeven. In een
populatie die nog klein is zal de drift ze binnen tientallen generaties
weer wegnemen, dus moet de redding worden herhaald of de populatie
worden vergroot.
\textbf{9.} $N_{e} = N$ bij gelijke geslachten: 500 volwassen dieren.
Vanaf 70 bij $r =
0.05$: $\ln(500/70)/0.05 = 39$ jaar.
\textbf{10.} De ideale populatie heeft Poisson-verdeelde
gezinsgrootten, met een variantie gelijk aan het gemiddelde van twee;
$N_{e} = 4N/(2 + V_{k})$, dus elk paar evenveel laten bijdragen
($V_{k} = 0$) geeft $N_{e} = 2N$: geen enkele lijn gaat door het geluk
van het nest verloren.
\textbf{11.} Zijn zaad verzamelen en bewaren voor kunstmatige
inseminatie, en het over de jaren bij verscheidene niet-verwante
vrouwtjes gebruiken --- de kosten zijn dat zijn allelen een kleine poel
zouden kunnen gaan overheersen, zodat zijn aandeel nabij $1/N_{e}$ moet
worden begrensd; of weefsel invriezen voor later klonen of het maken
van gameten, tegen de kosten van uitstel en technisch risico. Niets
doen verliest de allelen van de lijn voorgoed.
\textbf{12.} \qty{160}{} miljoen dollar voor zo'n 430 vogels: ongeveer
\num{370000} per vogel. Duur per papegaai, goedkoop tegen een soort
afgezet --- en minder dan een paar kilometer snelweg, wat de
vergelijking is die begrotingen werkelijk maken.
\textbf{13.} $(300/2000)^{0.25} = 0.62$: ongeveer 112 van de 180
soorten behouden, een uitstervingsschuld van 68.
\textbf{14.} Elke snipper $(100/2000)^{0.25}\times 180 = 85$ soorten.
Geheel verschillend: tot 255, onmogelijk boven de voorraad van 180;
gelijk: 85. Wat beslist is hoe verschillend de leefgebieden van de
snippers zijn en hoeveel soorten in \qty{100}{km^2} stand kunnen
houden; realistisch tussen 85 en 112.
\textbf{15.} $(320/2000)^{0.25}\times 180 = 114$, tegen 85 voor drie
geïsoleerde gelijke snippers: de corridor laat herkolonisatie en
genenstroom drie kleine populaties zich als één grotere gedragen.
\textbf{16.} $N_{e} = 50$ bij gelijke geslachten is 25 paren: \qty{1250}{km^2}.
De hele \qty{300}{km^2} herbergt zes paren, $N_{e} \approx 12$; een
snipper twee. Geen van beide volstaat: het roofdier gaat verloren wat
het aantal soorten ook is, en de oppervlakte beschermen die het nodig
heeft beschermt alles wat kleiner is --- een parapluesoort.
\textbf{17.} $\ln(200/10)/0.03 = 100$ jaar: de schuld wordt over een
eeuw betaald, lang nadat het bos is gekapt.
\textbf{18.} Randen nemen elke snipper zijn binnenkant af, zodat de
bruikbare oppervlakte kleiner is dan de gekarteerde en het verlies
groter dan de wet voorspelt; een matrix van secundair bos of van heggen
die sommige soorten kunnen gebruiken voegt effectieve oppervlakte toe
en maakt het verlies kleiner.
\textbf{19.} $r = \ln(100/31)/8 = \qty{0.146}{yr^{-1}}$;
verdubbelingstijd $\ln 2/0.146 = 4.7$ jaar.
\textbf{20.} $N(8) = 120/[1 + (120/31 - 1)\mathrm{e}^{-1.17}] = 120/(1 +
2.87\times 0.31) = 63$. De waargenomen 100 is de exponentiële waarde
zelf --- $r$ is juist op deze twee punten gepast --- tegenover de
logistische 63: in de eerste jaren stuitten de wolven op geen
dichtheidsafhankelijkheid --- lege territoria en overvloedige wapiti's.
\textbf{21.} $\ln(4000/\num{17000})/20 = -0.072$: \qty{7}{\%} per jaar.
Honderd wolven die \num{2000} wapiti's per jaar uit een kudde van
gemiddeld \num{8000} nemen is een kwart per jaar --- ruim genoeg om de
achteruitgang op eigen kracht te sturen, al deelden de jacht buiten het
park, de droogte en de predatie van kalveren door grizzlyberen erin.
\textbf{22.} Zestig volwassen exemplaren, onder 250: Ernstig Bedreigd.
Het roofdier bij $N_{e} \approx 12$ (een paar dozijn dieren): eveneens
Ernstig Bedreigd.
\textbf{23.} Van $3\times 10^{9}$ naar $1$ in 44 jaar: $r = -\ln(3\times
10^{9})/44 = -0.50$ per jaar, een halvering elke zeventien maanden.
Geen enkele oogst neemt elk jaar de helft van miljarden weg; de daling
verliep eerst trager en aan het eind veel sneller, omdat een koloniale
broeder onder een drempeldichtheid ophoudt zich voort te planten
terwijl zijn roofdieren niet langer worden verzadigd --- het
\omterm{def:b3:conservation-biology:small}{Allee-effect} maakte van een achteruitgang een ineenstorting.
\textbf{24.} In 1850 telde de populatie miljarden, was haar gemeten
groeitempo positief en haar variantie klein, en een \omterm{met:b3:conservation-biology:pva}{analyse van de
levensvatbaarheid} extrapoleert de dynamiek die zij krijgt: zij zou het
risico op nul hebben gesteld. Zij zou de aandrijving van buiten hebben
gemist --- de spoorwegen en de telegraaf waarmee jagers elke
slaapplaats konden volgen en leeghalen, het rooien van bos --- en de
Allee-drempel in een koloniale broeder, die geen van beide in de
demografie van een overvloedige soort zitten; zij extrapoleert, en de
oorzaken van uitsterven zijn gewoonlijk nieuw.
\textbf{25.} $N_{e} \approx 53$ vandaag; $F \approx 0.12$ na tien
generaties meer bij 60 (\qty{7}{\%} van de fitness); het gekrompen bos
behoudt ongeveer \qty{62}{\%} van zijn soorten.
\end{solution}
